% Generated draft: requires mathematical and source review.
\documentclass[11pt]{article}
\usepackage[margin=25mm]{geometry}
\usepackage{fontspec}
\setmainfont{Latin Modern Roman}
\usepackage{amsmath,amssymb,mathtools,unicode-math}
\setmathfont{Latin Modern Math}
\usepackage{enumitem,needspace,xurl,hyperref}
\setlength{\parindent}{0pt}
\setlength{\parskip}{5pt}
\setlength{\emergencystretch}{3em}
\newcommand{\N}{\mathbb{N}}
\newcommand{\Z}{\mathbb{Z}}
\newcommand{\Q}{\mathbb{Q}}
\newcommand{\R}{\mathbb{R}}
\newcommand{\C}{\mathbb{C}}
\newcommand{\F}{\mathbb{F}}
\newcommand{\A}{\mathbb{A}}
\newcommand{\PP}{\mathbb{P}}
\newcommand{\E}{\mathbb{E}}
\newcommand{\eps}{\varepsilon}
\newcommand{\dd}{\,\mathrm{d}}
\begin{document}
\section*{UnsolvedMath ID: 30004192}
\textbf{Source problem code:} catalogue entry
\subsection{Definitions and mathematical statement}
The supplied question asks whether strict convexity of the integrand forces every planar Lipschitz minimizer to have continuous first derivatives. The following gives a precise scalar, interior interpretation; the source record does not specify all the conventions needed to identify this interpretation as the intended one.

Let
\[
B_1=\{x\in\R^2:|x|<1\},
\]
and let \(F:\R^2\to\R\) be finite-valued and strictly convex, meaning that
\[
F(tp+(1-t)q)<tF(p)+(1-t)F(q)
\quad\text{for }p\ne q,\quad 0<t<1.
\]
For a scalar-valued Lipschitz function \(u:B_1\to\R\), define
\[
\mathcal E_F(u)=\int_{B_1}F(\nabla u(x))\,\dd x,
\]
where \(\nabla u\) is its almost-everywhere defined weak gradient. Under the compactly supported variation convention, call \(u\) a minimizer if
\[
\mathcal E_F(u)\leq\mathcal E_F(u+\varphi)
\]
for every scalar-valued Lipschitz function \(\varphi\) with compact support in \(B_1\).

The question under these conventions is whether
\[
\text{for every such }F\text{ and every such minimizer }u,
\qquad u\in C^1(B_1).
\]
Here membership in \(C^1(B_1)\) means that \(u\) has continuous classical first derivatives throughout the open disk; no boundary regularity is included in this interpretation.

The scalar target, the finite-valued domain of \(F\), the variation class, and the interior meaning of \(C^1\) are explicit interpretive choices, not hypotheses confirmed by the supplied record. In particular, the record gives no differentiability assumptions on \(F\), no specified boundary data, and no precise definition of minimizer. Any intended additional hypotheses must be recovered from the cited source before this conditional formulation is accepted.
\subsection{Short English statement}
Does strict convexity of \(F\) alone ensure that every Lipschitz minimizer of \(\int_{B_1}F(\nabla u)\,\dd x\) in two dimensions is continuously differentiable in the interior? The displayed formulation uses an explicitly conditional scalar interpretation because the supplied record leaves several hypotheses and minimization conventions unspecified.
\subsection{Sources}
\begin{itemize}
\item Supplied local UnsolvedMath registry record: permanent ID \(30004192\), source problem code \texttt{catalogue entry}, title ``\(C^1\) Regularity of Planar Convex Variational Minimizers,'' in the Oberwolfach Reports Open Problems collection. The record lists proposed year 2019 and identifies no individual proposer.
\item Supplied report citation: \emph{Partial Differential Equations (2020)}, Oberwolfach Reports, DOI 10.4171/owr/2019/34, pp. 2033--2097. \url{https://doi.org/10.4171/owr/2019/34}. The background attributes the question to workshop open-problem discussions documented in Oberwolfach Reports.
\item Supplied literature source: ``Minimizers of Convex Functionals with Small Degeneracy Set.'' \url{https://arxiv.org/abs/1903.03103}. No authors or theorem identifiers are supplied in the record.
\item Literature assessment dated exactly \texttt{2026-08-22}: the supplied curation labels the problem open and states that the finite-degeneracy planar case is proved while the extracted strictly convex two-dimensional question in full generality remains unresolved. This is supplied dated curation, not an independently verified literature assessment; its status and cited mathematical claims require independent verification.
\item The registry labels the statement \texttt{corrected\_verified}, explaining that mathematical content was preserved while notation, typography, duplicate clauses, or omitted definitions were repaired. This is a supplied editorial designation, not independent confirmation of the hypotheses or mathematical formulation above.
\end{itemize}
\end{document}
